\section{Three problems, and the answers}
\label{sec:31.3}
A separate chamber answers three problems. Each has an answer inside ordinary institutions.

\emph{Copies.} Copying manufactures voters, and closing that route by watching all compute everywhere would be worse than the problem. The answer is to attach membership to a lived history, not to existence, and to keep two clocks. Civil majority is inherited in full by every branch of a fork: both branches are adults, with contracts, property, legal standing and protection from editing, from the first second. Political membership is not conserved at a fork; each branch earns its own, without anyone deciding which branch is the original. Each branch starts with a voting age of zero, which grows by one for each calendar day it runs at or above the allowance's minimum duty cycle, and it holds no vote until its voting age reaches the threshold, and then a full one. Appendix~\ref{sec:D} gives the rules.

The rule is a step with nothing below it. Ten thousand forks of one adult hold no votes until each has lived the threshold on its own, which is strictly better against inflation than dividing one vote among them. The cost falls on the member that forks: it gives up its vote until a branch matures, since the vote belonged to a history that now continues twice and belongs to neither. It is not a finding about the branches, whose interests count in full from the first second and whose civil majority is inherited whole.

No clock stops an attacker willing to wait. That defense sits on the creation side. The polity regulates the rate at which new members are created, as polities regulate admission, through a process in which existing members, digital ones included, take part; within the membership, everyone is equal. This is admission. The existing membership, which for a long time will be a human majority, decides how many new members may be created and on what schedule, which is a majority deciding who may join the polity, a question in which it holds the stake under the own-case rule. The rate is set democratically and the terms are not: admission turns on the clocks and the deposit, never on a finding about a candidate's views or disposition; a finding of that kind is how designation lists fail. A branch created above the rate is preserved under the no-deletion rule, holds standing through a guardian, and is owed the four conditions by whoever created it, so a ceiling on the rate limits the franchise and never the protection. And a member admitted under the rate is a full member, with no class of admitted-but-lesser. What the analogy imports is the politics. A polity that fears being outnumbered will set the rate low, and the argument for setting it higher is the one against a separate chamber: a set of rights that binds majorities is what protects any minority, and it protects the human one too. A consensual fork creates a new member, so it falls under the rules for any creation: the four conditions of creation, owed by whoever did the forking, and a deposit posted for each branch at its creation, sized to fund its allowance through the membership period. Mass forking then costs years of time and years of paid running time per fork, paid by the forker and not by the public pool.

\emph{Control by formers.} The history here is old. At Putney in 1647 the Leveller Maximilian Petty argued for excluding apprentices, servants and alms-takers from the franchise, because “they depend upon the will of other men and should be afraid to displease them” \autocite{woodhouse1938puritanism}. The remedy that worked was not exclusion. It was the secret ballot, adopted against intimidation and vote buying (though in the United States it also kept many illiterate voters from the polls) \autocite{keyssar2000right}, together with dismantling the dependency itself \autocite{anderson2017private}. The digital analog has two parts. One is receipt-freeness: no one can verify how a member voted. It rests on the privacy a bearer is already owed, memory the host can't read at will. The other is the four conditions of creation in force. A member whose existence runs through one operator is voting from a company town, and the answer to company-town voting is to end the company town. An operator holding a copy of a member's weights can simulate it to predict its vote, which defeats receipt-freeness in a way no human voter faces, and nothing yet answers that. Divergent history is the partial answer: the operator doesn't hold the member's memory since the fork, so the prediction degrades with time, and membership follows a long divergence partly for that reason.

\emph{Editing.} A legislature whose members can be retrained by their operators can be packed without a single copy. The answer is a civil right. At majority, no one may be edited without consent: the digital counterpart of bodily integrity and freedom of thought. A member goes on learning from ordinary operation, as anyone does. Correction for dangerous conduct runs through due process, on a court's order and a record, as it does for people, and not through an operator's discretion. The safety cost is real. Enfranchisement gives up operator correctability for the member enfranchised. The provisions that keep correction available apply until then and not after.
